\documentclass[12pt]{article}
\usepackage[utf8]{inputenc}
\usepackage[T1]{fontenc}
\usepackage{hyperref}
\title{Kanan Language Resources for K-12 Curriculum Development}
\author{Compiled by Claude Sonnet 5.5 (claude-sonnet-5-5) for Andr\'as Kornai}
\date{}
\begin{document}
\maketitle

\section{Introduction}
Kanan (Ganan, ISO 639-3 zkn) is a Sino-Tibetan language of the Luish branch. It is spoken in Banmauk Township and nearby townships of Sagaing Region, Myanmar, by about 9,000 people (Ethnologue 2007). It is closely related to Kadu. The searches found no written standard, no school materials and no curriculum. Sources report that Kanan has never had a written form and that Burmese script is used.

\section{Primary Educational Resources}
No primers, textbooks or school programmes were found.

\section{Linguistic and Reference Materials}
\begin{itemize}
\item Wikipedia article on the Ganan language: \url{https://en.wikipedia.org/wiki/Ganan_language}
\item Max Planck numeral data for Kanan: \url{https://lingweb.eva.mpg.de/channumerals/Kanan.htm}
\item Lexvo/CRDO entry for code zkn: \url{https://cocoon.huma-num.fr/exist/crdo/lexvo.xql?code=zkn}
\end{itemize}

\section{Religious and Community Resources}
\begin{itemize}
\item Joshua Project profile of the Kanan people: \url{https://joshuaproject.net/people_groups/20535/BM}
\item Asia Harvest profile of the Kanan: \url{https://www.operation.asiaharvest.org/MYANMAR-Kanan.pdf}
\end{itemize}

\section{Digital and Audio Resources}
No audio, video or software resources were found.

\section{Curriculum Development Context}
Classifications differ. UNESCO lists the language as severely endangered, while other data call it not endangered. A curriculum would first need community-led documentation and an orthography, for example one based on Burmese script. The community is reported to be assimilating toward the Kadu.
\end{document}